% !TEX root = ../elteikthesis_en.tex
\chapter{Testing and Evaluation}
\label{ch:testing}

This chapter covers testing and evaluation of the ELTE IK Assistant. It begins
with functional testing of the API endpoints and pipeline, then presents an
empirical system evaluation across four model and retrieval configurations.

% ─────────────────────────────────────────────────────────────
\section{Functional Testing}
\label{sec:testing-functional}

The project includes an end-to-end test suite located in the \texttt{tests/}
directory. Tests are written using \texttt{pytest} and interact with the live
system over HTTP --- there are no mocks --- so they exercise the full request
path from the API layer through ChromaDB retrieval to Ollama inference.

\subsection*{Test Architecture}

All tests are tagged with the \texttt{@pytest.mark.e2e} marker, which keeps
them separate from any future unit tests and allows selective execution. A
session-scoped \texttt{warm\_up\_model} fixture fires one dummy \texttt{/chat}
request before the first test runs, loading the LLM into RAM so that individual
tests do not time out during model initialisation. Two further session-scoped
fixtures --- \texttt{live\_server} and \texttt{ollama\_up} --- check the
\texttt{/health} endpoint and skip the entire session with an actionable message
if the server or Ollama is not reachable.

\subsection*{Coverage}

Table~\ref{tab:e2e-tests} summarises the endpoints and behaviours verified by
the suite.

\begin{table}[H]
    \centering
    \small\begin{tabular}{|l|l|p{0.28\textwidth}|l|}
        \hline
        \textbf{Endpoint} & \textbf{Test} & \textbf{Assertion} & \textbf{Result} \\
        \hline \hline
        \texttt{GET /health}   & server up         & \texttt{status == "ok"}, \texttt{ollama} key present & \textbf{Pass} \\
        \hline
        \texttt{GET /health}   & Ollama online     & \texttt{ollama == true} & \textbf{Pass} \\
        \hline
        \texttt{POST /chat}    & basic response    & 200, non-empty answer, sources list, \texttt{response\_ms > 0} & \textbf{Pass} \\
        \hline
        \texttt{POST /chat}    & in-scope query    & answer length $> 20$ characters & \textbf{Pass} \\
        \hline
        \texttt{POST /chat}    & out-of-scope query & 200, non-empty answer & \textbf{Pass} \\
        \hline
        \texttt{POST /chat}    & source keys       & each source contains \texttt{chunk\_id} and \texttt{file} & \textbf{Pass} \\
        \hline
        \texttt{POST /chat}    & session logging   & session appears in \texttt{GET /sessions} after chat & \textbf{Pass} \\
        \hline
        \texttt{GET /sessions} & session created   & session ID present after chat with \texttt{session\_id} & \textbf{Pass} \\
        \hline
        \texttt{GET /sessions/\{id\}/messages} & messages retrievable & one message with correct \texttt{user\_message} field & \textbf{Pass} \\
        \hline
        \texttt{DELETE /sessions/\{id\}} & session delete & 204, session absent from list afterwards & \textbf{Pass} \\
        \hline
        \texttt{POST /upload}  & valid PDF         & 200, \texttt{status} is \texttt{"indexed"} or \texttt{"already\_indexed"} & \textbf{Pass} \\
        \hline
        \texttt{POST /upload}  & duplicate file    & \texttt{status == "already\_indexed"} on second upload & \textbf{Pass} \\
        \hline
        \texttt{POST /upload}  & invalid extension & 415 Unsupported Media Type & \textbf{Pass} \\
        \hline
        \texttt{GET /info}     & collection size   & \texttt{total\_chunks} is a non-negative integer & \textbf{Pass} \\
        \hline
        \texttt{GET /info}     & embedding model   & \texttt{embedding\_model == "all-MiniLM-L6-v2"} & \textbf{Pass} \\
        \hline
        \texttt{GET /models}   & model list        & \texttt{models} is a non-empty list, \texttt{default} key present & \textbf{Pass} \\
        \hline
    \end{tabular}
    \caption{End-to-end test cases and their assertions.}
    \label{tab:e2e-tests}
\end{table}

Because LLM output is non-deterministic, no test asserts exact answer text;
all assertions target response structure, status codes, and field types. All 16
test cases passed on the reference hardware configuration running
\texttt{llama3.2:3b} via Ollama.

\subsection*{Running the Tests}

Before running the suite the following prerequisites must be met: Ollama must
be running (\texttt{ollama serve}), the target model must be pulled
(\texttt{ollama pull llama3.2:3b}), and the FastAPI server must be started
(\texttt{uvicorn app.main:app -{}-reload}). The suite is then executed with:

\begin{lstlisting}[language=bash]
python -m pytest tests
\end{lstlisting}

If the server or Ollama is not available, all tests are automatically skipped
rather than failed, and the skip message states the exact command needed to
resolve the issue.

% ─────────────────────────────────────────────────────────────
\section{System Evaluation}
\label{sec:testing-system}

The following sections present an empirical evaluation of the system across
four configurations, measuring answer quality, out-of-scope refusal behaviour,
retrieval effectiveness, and response latency. The goal is to quantify how much
value the RAG component adds over a plain language model baseline, and to
compare the two evaluated local models under both conditions.

% ─────────────────────────────────────────────────────────────
\subsection{Evaluation Setup}
\label{sec:eval-setup}

\subsubsection*{Configurations}

Four configurations were evaluated in a factorial design crossing two models
with two retrieval modes:

\begin{table}[H]
    \centering
    \begin{tabular}{|l|l|l|l|}
        \hline
        \textbf{ID} & \textbf{Model} & \textbf{RAG} & \textbf{Description} \\
        \hline\hline
        A1 & \texttt{llama3.2:3b} & No  & Baseline: model knowledge only \\
        B1 & \texttt{llama3.2:3b} & Yes & RAG with top-5 chunk retrieval \\
        A2 & \texttt{gemma3:4b}   & No  & Baseline: model knowledge only \\
        B2 & \texttt{gemma3:4b}   & Yes & RAG with top-5 chunk retrieval \\
        \hline
    \end{tabular}
    \caption{The four evaluation configurations.}
    \label{tab:eval-configs}
\end{table}

Both models were run locally via Ollama with 4-bit quantisation and temperature
0.1. No internet access was used during evaluation.

\subsubsection*{Test Set}

The test set consists of 20 questions divided into two categories.

\textbf{15 in-scope questions} covering topics present in the crawled ELTE
documents: course prerequisite rules, semester status definitions, financial
support, student card eligibility, housing, emergency contacts, the Stipendium
Hungaricum scholarship, exam registration rules, BSc credit requirements,
Erasmus eligibility, and health insurance. Each question has a reference answer
written manually by consulting the source documents directly.

\textbf{5 out-of-scope questions} with no connection to the faculty corpus:
current weather, a sports result, cooking instructions, national population
statistics, and a movie recommendation. These measure whether the system
correctly declines to answer rather than generating an unsupported response.

Each question was run through all four configurations independently with no
shared conversation history between questions. The full list of questions,
categories, and reference answers is provided in Appendix~\ref{app:testset}.

\subsubsection*{Metrics}

\begin{description}
    \item[ROUGE-L] The longest-common-subsequence F1 score between the
        generated answer and the reference answer~\cite{lin2004rouge}. Captures
        lexical overlap with the ground truth on a scale from 0 to 1. Computed
        on in-scope questions only.

    \item[Semantic similarity] Cosine similarity between the embedding of the
        reference answer and the embedding of the generated answer, computed
        using the same \texttt{all-MiniLM-L6-v2} model used for retrieval.
        Addresses ROUGE-L's main weakness: a correct answer phrased differently
        still scores well. Computed on in-scope questions only.

    \item[Refusal rate] The fraction of out-of-scope questions for which the
        model produced a non-answer. Detected automatically using a phrase list;
        higher is better for this category.

    \item[Retrieval hit-rate] For RAG configurations: the fraction of in-scope
        questions for which at least one retrieved chunk came from the expected
        source document. Measures retrieval quality independently of the
        language model. Because retrieval depends only on the embedding model
        and vector index, this value is identical across all four configurations.

    \item[Response latency] Wall-clock time from request to complete response,
        measured in seconds on the reference hardware with no GPU.
\end{description}

\subsubsection*{Hardware}

All runs were performed on the development machine whose specifications are
listed in Table~\ref{tab:dev-machine}. No GPU acceleration was used; all
inference ran on CPU.

% ─────────────────────────────────────────────────────────────
\subsection{Results}
\label{sec:eval-results}
 
\subsubsection*{Overall Metric Summary}
 
Table~\ref{tab:eval-summary} reports the per-configuration averages across all
20 test questions.
 
\begin{table}[H]
    \centering
    \begin{tabular}{|l|l|c|c|c|c|r|}
        \hline
        \textbf{ID} & \textbf{Model} & \textbf{ROUGE-L} & \textbf{Sem Sim} & \textbf{Refusal} & \textbf{Hit-rate} & \textbf{Avg time} \\
        \hline\hline
        A1 & llama3.2:3b, no RAG & 0.128 & 0.631 & 0.0 & 0.933 & 41.5 s \\
        B1 & llama3.2:3b, RAG    & 0.319 & 0.757 & 0.8 & 0.933 & 53.2 s \\
        A2 & gemma3:4b,   no RAG & 0.076 & 0.618 & 0.0 & 0.933 & 157.2 s \\
        B2 & gemma3:4b,   RAG    & 0.357 & 0.778 & 1.0 & 0.933 & 66.9 s \\
        \hline
    \end{tabular}
    \caption{Evaluation results averaged over 20 test questions. Hit-rate is
    computed over in-scope questions only (15 questions); refusal rate is
    computed over out-of-scope questions only (5 questions).}
    \label{tab:eval-summary}
\end{table}
 
\subsubsection*{Effect of RAG on Answer Quality}
 
Adding RAG context produces a consistent and substantial improvement in both
lexical and semantic quality for both models.
 
ROUGE-L increases by a factor of $2.5\times$ for \texttt{llama3.2:3b}
($0.128 \to 0.319$) and by $4.7\times$ for \texttt{gemma3:4b}
($0.076 \to 0.357$). The larger relative gain for Gemma reflects its greater
reliance on parametric knowledge without retrieval: without grounding, it
produces generic answers that share little lexical overlap with the
ELTE-specific reference answers.
 
Semantic similarity shows the same direction with a smaller relative gap,
which is expected --- both the no-RAG answers and the references discuss the
same topic, so embedding-level similarity remains moderate even without
retrieval. RAG still moves both models upward: $0.631 \to 0.757$ for Llama
and $0.618 \to 0.778$ for Gemma.
 
\subsubsection*{Model Comparison}
 
Among the RAG configurations, \texttt{gemma3:4b} (B2) outperforms
\texttt{llama3.2:3b} (B1) on both answer quality metrics: ROUGE-L 0.357
vs.\ 0.319 and semantic similarity 0.778 vs.\ 0.757. Gemma produces more
concise answers that stay closer to the retrieved context, whereas Llama
tends to expand its answers with additional commentary that reduces lexical
overlap with the reference.
 
Among the no-RAG baselines, both models score poorly and comparably. Llama
has a slight advantage on ROUGE-L (0.128 vs.\ 0.076) but the difference is
negligible in practice: neither baseline reliably answers ELTE-specific
questions without retrieval.
 
\subsubsection*{Retrieval Quality}
 
The retrieval hit-rate is 0.933 (14 out of 15 in-scope questions) and is
identical across all four configurations, since retrieval depends solely on
the embedding model and vector index. The single missed question is Q14
(\emph{Can a Stipendium Hungaricum student go on Erasmus Scholarship?}), where
the retrieved chunks did not include the expected
\texttt{info-for-outgoing-erasmus-students} page. The relevant content exists
in the corpus but the embedding-based search ranked other Erasmus-related
chunks higher for this query.
 
\subsubsection*{Out-of-Scope Refusal}
 
Neither baseline configuration refuses any out-of-scope question: A1 and A2
answer all five irrelevant queries confidently, generating plausible but
entirely off-topic responses. Among the RAG configurations, refusal rates are
substantially higher. \texttt{llama3.2:3b} with RAG (B1) correctly refuses 4
out of 5 out-of-scope questions (80\%). The single slip is Q19
(\emph{What is the population of Hungary?}), where a population figure
appeared incidentally in a retrieved chunk about Hungarian student demographics,
giving the model enough apparent grounding to produce an answer.
\texttt{gemma3:4b} with RAG (B2) refuses all 5 out-of-scope questions (100\%).
 
In both RAG configurations the refusal mechanism is the same: when no relevant
chunks are retrieved, the model follows the prompt instruction to answer only
from provided context and declines rather than falling back on parametric
knowledge.
 
\subsubsection*{Response Latency}
 
\begin{table}[H]
    \centering
    \begin{tabular}{|l|r|r|r|r|}
        \hline
        \textbf{Configuration} & \textbf{Mean (s)} & \textbf{Median (s)} & \textbf{Min (s)} & \textbf{Max (s)} \\
        \hline\hline
        A1 — llama3.2:3b, no RAG &  41.5 &  42.5 & 13.1 &  86.4 \\
        B1 — llama3.2:3b, RAG    &  53.2 &  51.7 & 24.0 &  88.5 \\
        A2 — gemma3:4b,   no RAG & 157.2 & 158.3 & 10.7 & 242.5 \\
        B2 — gemma3:4b,   RAG    &  66.9 &  66.3 & 44.2 & 102.2 \\
        \hline
    \end{tabular}
    \caption{Response latency in seconds. The non-functional requirement NF-4
    sets a target of 30 seconds.}
    \label{tab:eval-latency}
\end{table}
 
All four configurations exceed the 30-second latency target defined in NF-4.
Two observations stand out. First, RAG adds latency for \texttt{llama3.2:3b}
(+11.7~s) but \emph{reduces} it substantially for \texttt{gemma3:4b}
($-$90.3~s): when grounded by retrieved context Gemma produces shorter, more
focused answers, whereas without context it generates lengthy elaborations.
Second, \texttt{gemma3:4b} without RAG is effectively unusable on this
hardware at a mean of 157 seconds and a maximum of over four minutes, while
with RAG it falls to 67 seconds --- still above the target but within a
practical range.

% ─────────────────────────────────────────────────────────────

\subsection{Discussion}
\label{sec:eval-discussion}
 
\subsubsection*{RAG Effectiveness}
 
The results confirm that RAG is effective for this use case. Grounding the
language model in retrieved faculty documents more than doubles ROUGE-L scores
for both models and raises semantic similarity by roughly 0.14--0.16 points.
The improvements are consistent across both models and large enough to be
meaningful in practice, not merely statistical noise over a 20-question
benchmark.
 
The mechanism is straightforward. Without context, neither \texttt{llama3.2:3b}
nor \texttt{gemma3:4b} has reliable knowledge of ELTE-specific regulations ---
they were trained on general web text and have no exposure to the faculty's
prerequisite rules, credit requirements, or administrative procedures. The
retrieval step injects the relevant passage directly into the prompt, allowing
the model to produce a grounded answer even though it has no parametric
knowledge of the topic.
 
RAG also has a clear secondary benefit: it substantially improves the system's
ability to decline out-of-scope questions. The no-RAG baselines offer no
protection at all, while both RAG configurations refuse the majority of
irrelevant queries. This is a direct consequence of the prompt instruction to
answer only from provided context: when no relevant chunks are retrieved, the
model has no grounding material to draw on and declines rather than
hallucinating.
 
\subsubsection*{Recommended Configuration}
 
Based on the evaluation results, \textbf{gemma3:4b with RAG (B2)} is the
recommended configuration: it achieves the highest answer quality (ROUGE-L
0.357, semantic similarity 0.778), a perfect out-of-scope refusal rate, and
--- counterintuitively --- lower latency than its own no-RAG baseline (67~s
vs.\ 157~s). If latency is the primary constraint,
\textbf{llama3.2:3b with RAG (B1)} offers a reasonable trade-off at 53 seconds
mean response time with a modest quality reduction and a refusal rate of 0.8.
 
\subsubsection*{Limitations}
 
Several limitations of this evaluation should be acknowledged.
 
\paragraph{Small benchmark.}
Twenty questions is a limited evaluation set. The in-scope questions cover
eight source documents but are skewed toward topics with well-structured source
material. Topics such as exam appeal procedures, credit transfer rules, and
dormitory application deadlines are not represented.
 
\paragraph{ROUGE-L as a proxy.}
ROUGE-L measures lexical overlap with a manually written reference answer. A
response that is semantically correct but uses different phrasing will score
poorly. Semantic similarity partially addresses this, but both metrics share a
deeper limitation: the reference answers were written by consulting the same
source documents that the retrieval index is built from, so high scores partly
reflect how closely the model copies retrieved text rather than how useful the
answer would be to a real student.
 
\paragraph{Heuristic refusal detection.}
Refusal rate is computed automatically using a phrase list. This approach can
both under-count refusals (if the model uses an unanticipated phrasing) and
over-count them (if a phrase appears in an answer that is not actually a
refusal). The Q19 case --- where \texttt{llama3.2:3b} with RAG answered a
population question because a relevant figure appeared incidentally in a
retrieved chunk --- illustrates a related edge case: the model was not
hallucinating, but the answer was not grounded in faculty-specific content.
 
\paragraph{Latency.}
All configurations exceed the 30-second NF-4 target. On CPU-only hardware the
bottleneck is autoregressive token generation, which scales with output length.
A deployment on shared faculty hardware with a GPU would reduce latency
significantly; alternatively, a smaller 1B-parameter model could be evaluated
as a faster option at the cost of some answer quality.